\section{Methodology}

The proposed system processes English loanwords through a pipeline of four
sequential stages: graphemic tokenization, rule-based grapheme mapping,
phonetic vowel resolution, and output normalization. Each stage operates on the
output of the previous, and together they model the multi-layered nature of
loanword adaptation.

\subsection{Graphemic Tokenization}

Before any mapping can take place, the English input must be segmented into its
functional graphemic units. This step is necessary because English orthography
contains multi-letter sequences that function as single phonological units. The
system recognizes the following classes of multigraphs:

\begin{itemize}
    \item \textbf{Digraphs}: \emph{ph}, \emph{ch}, \emph{sh}, \emph{th},
          \emph{ng}, \emph{kn}, \emph{wr}, \emph{gh}, \emph{ee}, \emph{oo}
    \item \textbf{Trigraphs}: \emph{igh}, \emph{tch}, \emph{dge},
          \emph{que}, \emph{ore}
    \item \textbf{Higher-order multigraphs}: \emph{augh}, \emph{ough},
          \emph{eigh}
\end{itemize}

For instance, \emph{ch} in ``chapter'' represents a single phonemic unit, the
voiceless postalveolar affricate /\textipa{tS}/, and must be treated as one
token rather than two independent letters. Similarly, \emph{igh} in ``light''
encodes the diphthong /\textipa{aI}/ as a single unit.

The tokenizer applies a greedy, longest-match strategy: at each position in the
input string, it attempts to match a tetragraph first, then a trigraph, then a
digraph, and finally falls back to a single character if no multi-character
pattern is recognized. Formally, for each position $i$ in the input string, the
tokenizer selects:
\[
    \max \{n\in\{1,2,3,4\} \mid \mathrm{substr}(i,n) \in \Sigma\}
\]
where \(\Sigma\) is the set of recognized graphemic units and
\(\text{substr}(i, n)\) is the \(n\)-character substring beginning at position
\(i\). For example, ``straight'' and ``phishing'' are tokenized as:

\begin{center}
    \texttt{S~T~R~Aigh~T} \\
    \texttt{Ph~I~Sh~I~Ng}
\end{center}

This ensures that subsequent mapping rules operate on phonologically meaningful
units rather than on raw characters.

\subsection{Rule-Based Grapheme Mapping}

Following tokenization, each graphemic token is mapped to its Filipino
orthographic equivalent through a priority-ordered cascade of rule types. For
each token, the system tests the following rule classes in order, accepting the
first match:

\begin{enumerate}
    \item Context-sensitive orthographic rules
    \item Phonetic rules
    \item Context-free orthographic rules
\end{enumerate}

This ordering reflects a key design principle: more specific rules take
priority over more general ones. Context-sensitive rules that depend on
neighboring graphemes are checked first because they handle the majority of
ambiguous cases; context-free rules serve as the default fallback.

\subsubsection{Context-Free Rules}

Many English consonants bear an invariant relationship to their pronunciation
and therefore to their Filipino adapted form. For such graphemes, adaptation
reduces to a direct one-to-one substitution requiring no knowledge of
surrounding context. We call these \emph{context-free} rules. For example:

\begin{itemize}
    \item \(\emph{b} \mapsto \emph{b}\): both instances in ``basketball'' map
          directly, yielding \emph{basketbol}.
    \item \(\emph{ph} \mapsto (\emph{p}\mid\emph{f})\): the digraph \emph{ph}
          maps to \emph{p} or \emph{f},
          as in ``cellphone'' is adapted to ``selpon'' or ``selfon''.
\end{itemize}

Because these mappings are strictly positional and produce the same output
regardless of surrounding tokens, they are computed without reference to any
context window.

\subsubsection{Context-Sensitive Rules}

Many English graphemes have pronunciation that varies depending on surrounding
letters, and their Filipino adaptation must vary accordingly. We call these
\emph{context-sensitive} rules. The rule for a given token is selected by
inspecting adjacent graphemes and matching against a set of ordered conditions.
For example:

\begin{itemize}
    \item Cases for \emph{c}:
          \begin{enumerate}
              \item \(\emph{c} \cdot (\emph{i} \mid \emph{e}) \mapsto s\):
                    when \emph{c} precedes \emph{i} or \emph{e}, as in ``cinema'' \(\mapsto\) \emph{sine}.
              \item Otherwise: \emph{c} \( \mapsto \) \emph{k}, as in ``computer'' \(\mapsto\)
                    \emph{kompyuter}.
          \end{enumerate}
    \item Cases for doubled consonants: when two identical consonants appear adjacently
          (e.g., \emph{ll}, \emph{ss}), the second is treated as silent and mapped to no
          output, with the single consonant preserved.
\end{itemize}

More complex patterns require multi-token lookahead. The sequence \emph{-ate}
at the word boundary, for instance, maps to \emph{-eyt} rather than a sequence
of individual vowel and consonant adaptations, reflecting the holistic
treatment of common English suffixes in Filipino borrowing.

\subsection{Phonetic Resolution for Vowels}

The context-sensitive rules described above handle consonants well, but vowel
adaptation presents a qualitatively different challenge. Consider the grapheme
\emph{a} alone: it may represent /\textipa{\ae}/ in cat'', /\textipa{eI}/ in
cake'', /\textipa{A:}/ in father'', and /\textipa{@}/ in about''. Because
Filipino has only five vowels (/a/, /e/, /i/, /o/, /u/), the system must
collapse these distinctions; the correct collapsed form depends on the actual
pronunciation rather than the spelling. Furthermore, because English is a
stress-timed language, unstressed vowels frequently reduce to the schwa
/\textipa{@}/, introducing a systematic ambiguity that orthography alone cannot
resolve. For these tokens, the system invokes a G2P component to obtain an IPA
transcription of the entire word. The relevant vowel phoneme is then extracted
by aligning it to the current token.

\begin{itemize}
    \item Tense vowels and diphthongs (/\textipa{eI}/, /\textipa{oU}/, /\textipa{aI}/,
          /\textipa{aU}/) are rendered as two-character Filipino sequences: \emph{ey},
          \emph{ow}, \emph{ay}, \emph{aw}.
    \item Lax monophthongs (/\textipa{\ae}/, /\textipa{E}/, /\textipa{I}/, /\textipa{Q}/,
          /\textipa{U}/) are mapped to the nearest Filipino simple vowel.
    \item The schwa /\textipa{@}/ in word-final position is realized as \emph{e} or
          \emph{o} depending on the written vowel; in non-final position it is suppressed
          and the vowel is omitted from the output.
    \item R-colored vowels (/\textipa{@r}/, /\textipa{3r}/) preserve the written vowel
          letter, deferring to orthographic convention.
\end{itemize}

\subsection{Output Normalization}
After all per-token rules have been applied, the system performs a final
normalization pass over the raw output sequence. This pass corrects three
classes of artifact that can arise from the independent firing of rules across
adjacent tokens:
\begin{itemize}
    \item \textbf{Contiguous duplicates}: identical adjacent graphemes produced by separate rule branches are collapsed to a single instance (e.g., \emph{ll} \( \mapsto \) \emph{l}).
    \item \textbf{Redundant affricate prefixes}: when a stop consonant is immediately followed by its own affricate form (\emph{d} followed by \emph{dy}, or \emph{t} followed by \emph{ts}), the leading stop is removed, as it is already encoded within the affricate.
    \item \textbf{Final schwa-liquid clusters}: when the IPA transcription contains a schwa and the word-final two graphemes form an illegal consonant cluster ending in \emph{l} or \emph{r}, a vowel \emph{o} is inserted before the final liquid to restore syllabic well-formedness (e.g., \emph{-bl} \( \mapsto \) \emph{-bol}).
\end{itemize}

Together, these normalization steps ensure that the output conforms to the
phonotactic and orthographic conventions codified in the KWF \emph{Manwal sa
    Masinop na Pagsulat}~\cite{almario2014}.